\begin{claim}{C.38}{Proved}{Discrete cyclic operators on nonabelian line and circle orbits}
Let $B$ be a line or circle and let
$T\leq\operatorname{Homeo}^+(B)$ be transitive and
nonabelian. Its full compact-open centralizer
\[
                  N_B=C_{\operatorname{Homeo}(B)}(T)
                                                        \tag{C.38.1}
\]
is discrete and acts freely and properly. On a line it is
trivial or infinite cyclic; on a circle it is finite cyclic.

In the plane-loop setting of C.28, let $B$ be a nonopen
orbit whose intrinsic image $G|_B$ is nonabelian, and let
$H_B=\{h\in H:hB=B\}$ be the operator stabilizer.
Its actual restriction has the exact sequence
\[
 1\longrightarrow K_B\longrightarrow H_B
       \xrightarrow{\ \rho_B\ } P_B\longrightarrow1,
 \quad
 K_B=\{h\in H_B:h|_B=1\},\quad P_B=\rho_B(H_B)\leq N_B.
                                                        \tag{C.38.2}
\]
In the intrinsic compact-open topology, $P_B$ is already
closed and discrete, acts freely and properly, and
\[
                    p_B:B\longrightarrow B/P_B
                                                        \tag{C.38.3}
\]
is an actual regular covering. Its deck transformations
are exactly $P_B$. The quotient is a line for trivial
$P_B$ on a line, and otherwise a circle. On a circle
$P_B$ is finite and the covering degree is $|P_B|$.
The kernel $K_B$ acts identically on $B$ throughout this
comparison.
\end{claim}

\begin{proof}
Let $L$ be the closure of $T$ in the compact-open
homeomorphism group of $B$, choose $b\in B$, and put
$J=L_b$. The group $L$ is Polish and transitive, and
$C(L)=C(T)$ by continuity of composition.
The open orbit-map theorem used in C.26 gives
$L/J\cong B$. The normalizer $N_L(J)$ is closed, and
the same exact evaluation construction as in C.30 gives
\[
 \operatorname{ev}_b:N_B
       \xrightarrow{\ \cong\ }N_L(J)/J\ \subset B.
                                                        \tag{C.38.4}
\]
The right side is closed. The inverse sends the coset
represented by $k$ to the map $lb\mapsto lkb$.
The normalizer equation makes this well defined.
Its compact-open continuity follows by the same
microtransitivity argument as C.30: convergent points
of $B$ lift locally through the open orbit map to
elements of $L$ tending to the identity.
No surface-dimensional step occurs in that argument.
The construction also proves freeness, local compactness,
and properness. Explicitly, if
$n_i b_i=c_i$, $b_i\to b$, and $c_i\to c$, choose
$l_i\to1$ with $l_i b=b_i$. Then
$n_i b=l_i^{-1}c_i\to c$, and the closed image and
homeomorphism (C.38.4) give a convergent subsequence
in $N_B$. This is the required properness test.

Every orientation-reversing homeomorphism of a line
or circle has a fixed point. Freeness therefore places
$N_B$ in $\operatorname{Homeo}^+(B)$.
Its faithful product action on the three-manifold
$B\times\R^2$ makes it a Lie group by the verified
Hilbert--Smith theorem used in C.33.
Equation (C.38.4) bounds its dimension by one.
If that dimension were one, every orbit of $N_B^0$
would be open in $B$: its one-dimensional orbit map
is injective and hence open by invariance of domain.
Connectedness of $B$ would make $N_B^0$ transitive.
Freeness would then give $N_B=N_B^0$, since an element
of $N_B^0$ with the same evaluation as an element of
$N_B$ equals it. This is the translation group $\R$
or the rotation group $S^1$ in evaluation coordinates.
Its centralizer is that same abelian group of right
translations. But $T\leq C(N_B)$, contradicting the
assumed nonabelianness of $T$.
Consequently $N_B$ is a zero-dimensional Lie group,
so is discrete.

A free proper discrete action supplies evenly covered
neighborhoods, and its orbit space is a connected
one-manifold without boundary. For a line this
quotient is a line or a circle. Since the line is
simply connected, the first case has trivial deck
group and the second has deck group $\Z$.
For a circle the quotient is a circle and the
covering has finitely many sheets; its deck group
is finite cyclic. This proves the first assertion.

In the stated plane-loop application, C.28 proves
that each operator stabilizing $B$ induces an actual
homeomorphism for its intrinsic line or circle
topology. Commutation with $G$ puts its image in
(C.38.1). Every subgroup of this discrete group is
closed; restriction of a proper free action to a
closed subgroup is proper and free. These facts
give (C.38.3). Its fibers are precisely the
$P_B$-orbits. A deck transformation is determined
by one point in the connected covering, and each
point of a fiber is obtained by exactly one
element of $P_B$, so the whole deck group is
exactly $P_B$. The same one-manifold covering
classification gives the stated quotient and
degree. Equation (C.38.2) retains the original
stabilizer and its restriction kernel.
\end{proof}

\paragraph{Topology and the remaining relative extension.}
The topology asserted for $P_B$ is its intrinsic
compact-open image topology. The operator group may
be regarded as discrete as in the original problem;
no assertion about continuity of the restriction
map from an inherited plane topology to a recurrent
intrinsic line topology is needed.
C.38 gives exact covering quotients on every
nonabelian one-dimensional orbit. It supplies no
transverse control between different orbits and no
relative extension of the C.37 family from $X_2$
to its exterior. Those requirements remain in C.12.
